\scanpage{278}
\mspage{19}

pour les catégories réduites de longueur \uline{\unsure{pair|paire}}.

Comme chaque fois qu'on a un iso. $\theta$ d'int. \uline{orientés}, on trouve \deleted{cette fois} un iso. entre les catégories de chemins de type \deleted{$I$ et $I^{\mathrm{opp}}$} $(I,\omega)$ et $(I,\omega)^{\mathrm{opp}}=(I,\omega')$, \uline{compatible} cette fois avec les foncteurs $(\unsure{\sigma|\alpha},\beta)$ :
\diagram{\begin{tikzcd}[ampersand replacement=\&,column sep=small]
\underline{\mathrm{Ch}}_{(I,\omega)}(X) \arrow[rr,"{\substack{\unsure{\psi|\Psi}_{\theta}\\\sim}}"] \arrow[dr,"{(\sigma,\beta)}"] \&{} \& \underline{\mathrm{Ch}}_{(I,\omega')}(X) \arrow[dl,dash,"{(\sigma',\beta')}"']\\
{} \& X\times X \& {}
\end{tikzcd}}

\deleted{On peut regarder $\theta$ comme un \unsure{autom|autom.} non trivial de $I$ (échangeant \unsure{ses|les} deux extrémités) et qui opère \inserted{donc aussi} sur \deleted{\unsure{$\mathrm{Ch}$}} $\underline{\mathrm{Hom}}(I,X)$, et c'est ça \deleted{\unsure{la flèche}} l'iso. $\psi_\theta$ du \deleted{\unsure{le f}} diagramme précédent. Il transforme donc le \deleted{diagramme}\inserted{chemin de} type $\Phi_1$
\diagram{\begin{tikzcd}[ampersand replacement=\&,column sep=small,row sep=2pt]
{} \&{} \&{} \&{} \&{} \& a_0\\
a_0 \& a_1 \arrow[l] \arrow[r] \& a_2 \& \text{ou} \& a_1 \arrow[ur] \arrow[dr] \&{}\\
{} \&{} \&{} \&{} \&{} \&a_2
\end{tikzcd}}
\inserted{non pas en} en le chemin
\diagram{\begin{tikzcd}[ampersand replacement=\&,column sep=small,row sep=2pt]
{} \&{} \&{} \&{} \&{} \&a_2\\
a_2 \&a_1 \arrow[l] \arrow[r] \&a_0 \& \text{ou} \& a_1 \arrow[ur] \arrow[dr] \&{}\\
{} \&{} \&{} \&{} \&{} \&a_0
\end{tikzcd}}
\deleted{\unsure{Ils sont}} mais en lui-même. En tant que foncteur \deleted{\uline{Ces deux chemins} sont de même type $\Phi_1$, ($(I,\omega)\simeq(I,\omega')$) j'ai prétendu qu'il revient} \deleted{\unsure{à trouver|à trouve}} \inserted{sur $I$,} le chemin opposé à}
% archive sheet number: 277
